
\chapter{Discussions \& Conclusions}
\label{chap:discuss_and_conclus}



% \ifpdf
%     \graphicspath{{7_conclusions/figures/PNG/}{7_conclusions/figures/PDF/}{7_conclusions/figures/}}
% \else
%     \graphicspath{{7_conclusions/figures/EPS/}{7_conclusions/figures/}}
% \fi




{\color{red} 
Add a new Chapter 8: "Discussion and Conclusions" — approximately one page, free prose (no bullet
lists), structured as three paragraphs:
Paragraph 1: restate the research question and give the direct answer.
Paragraph 2: reflect on the difficulty of the project and the key learnings.
Paragraph 3: future work (moved from Chapter 7).
}

\section{Future Work}


\begin{itemize}
  \item \textbf{Repeat Experiment~3 with the intended 400K DQN buffer
  size.} This is the most direct follow-up: the buffer oversight means
  the Breakout result should be re-verified with the correct
  configuration.
  \item \textbf{Longer training budgets.} Running to 50M+ frames on
  Atari would bring the comparison closer to the original DQN and PPO
  paper settings, and may further change the per-environment ranking.
  \item \textbf{Larger $M$.} Adding more Atari environments
  (SpaceInvaders, MsPacman, BeamRider) and more classic-control
  environments (MountainCar, Pendulum) would make the cross-environment
  aggregates more stable and the conclusions more generalisable.
  \item \textbf{Hyperparameter tuning study.} Tuning one environment
  per algorithm would quantify the gap between ``default SB3'' and
  ``best achievable'', giving context to the default-hyperparameter
  numbers reported here.
  \item \textbf{Continuous-control extension.} A MuJoCo or
  continuous-control study would test whether the per-environment
  pattern seen here generalises beyond discrete-action domains, where
  policy-gradient methods are typically expected to have a larger
  advantage.
\end{itemize}


